\section{Materiales Auxeticos}
Los materiales y estructuras auxéticas se definen fundamentalmente por poseer un Coeficiente de Poisson Negativo (NPR) , . A diferencia de la gran mayoría de los materiales convencionales, que se adelgazan cuando se estiran y se ensanchan cuando se comprimen, los materiales auxéticos presentan el comportamiento contraintuitivo de expandirse transversalmente cuando se estiran y contraerse lateralmente cuando se comprimen , .
El término "auxético" fue acuñado en 1991 por el investigador Ken Evans y deriva de la palabra griega auxetikos, que significa "aquello que tiende a aumentar" o expandirse , . Esta inusual propiedad expansiva generalmente no deriva de la composición química del material en sí, sino de su geometría estructural interna y de sus mecanismos de deformación (tales como la flexión, rotación o el despliegue de sus componentes microscópicos) , . Gracias a este coeficiente negativo, los materiales auxéticos logran propiedades mecánicas excepcionales, como una altísima resistencia a la indentación, alta absorción de energía, gran tenacidad a la fractura y la capacidad de formar superficies en forma de cúpula (curvatura sinclástica) en lugar de forma de silla de montar al curvarlos , , .
Los diseños auxéticos pueden existir y aplicarse a cualquier escala, desde nivel macroscópico hasta el molecular, y se han desarrollado una amplia variedad de estructuras auxéticas, como los patrones de re-entrant (con forma de flecha), los patrones de chiral (con elementos giratorios) y los patrones de rotación (con elementos que rotan para expandirse) , . Estas estructuras se pueden fabricar utilizando técnicas avanzadas como la impresión 3D, la fabricación aditiva y la microfabricación , . Los materiales auxéticos han encontrado aplicaciones en campos tan diversos como la ingeniería civil, la aeroespacial, la biomedicina, la protección personal y el diseño de productos , , .
